\documentclass[10pt,a4paper]{amsart}
\usepackage[margin=19mm]{geometry}
\usepackage{amsmath,amssymb,mathtools}
\usepackage[scr=rsfs,cal=euler]{mathalfa}
\usepackage{xfrac}
\usepackage[unicode,hidelinks,bookmarks=false]{hyperref}
\newcommand{\ie}{i.e.\ }
\newcommand{\cf}{cf.\ }
\newcommand{\resp}{respectively\ }
\newcommand{\Subsection}{\subsection}
\providecommand{\nobreakdash}{\nobreak\hbox{-}}
\setlength{\emergencystretch}{2em}
\multlinegap0pt
\usepackage{amssymb}
\usepackage{xcolor}
\newtheorem{theorem}{Theorem}
\newcommand{\abs}[1]{\left|#1\right|}
\newcommand{\lam}{\lambda}
\newcommand{\pvec}{\vct{p}}
\newcommand{\qvec}{\vct{q}}
\newcommand{\vct}[1]{\mathbf{#1}}
\newcommand{\yvec}{\vct{y}}
\title{TV over Bernoulli products: the small parameter regime}
\author{Ariel Avital, Aryeh Kontorovich, George Salafatinos}
\date{Source: arXiv:2602.21828v1 (CC BY 4.0)}
\begin{document}
\maketitle

\noindent\emph{Source and licence.} TV over Bernoulli products: the small parameter regime. Version arXiv:2602.21828v1 (CC BY 4.0), distributed by arXiv under the Creative Commons Attribution 4.0 International licence, http://creativecommons.org/licenses/by/4.0/, as recorded in the arXiv licence metadata for this exact version and shown on the abstract page https://arxiv.org/abs/2602.21828v1. Full licence notice: \texttt{licenses/arxiv-2602.21828-CC-BY-4.0.txt}.\par\smallskip

\noindent\textit{Editorial scope.} Counted unit: the article's theorem thm:Delta2\_bound (source0.tex lines 474--479). The article's complete original proof of it is delivered with it (source0.tex lines 481--567, one \texttt{proof} environment of 87 lines). No mathematics is written, repaired or re-derived: the statement and the proof are the article's own lines, in the article's own notation, and no retained line is substituted or re-flowed.  Retained uncounted as the source-local context the proof needs: lines 397--403.  Nothing was omitted from any retained span, so the package carries no omission record.  No retained line contains a citation, so the wrapper carries no bibliography and no bibliography entry.  No retained line contains a figure environment, an image-inclusion command or drawing code, so no image is delivered and the package declares no asset dependency.  Editorial formatting only, in addition: an AMS \texttt{amsart} preamble that restates the article's own declarations and macros used by the retained lines, namely the environments \texttt{theorem} on separate counters, together with the standard packages those lines need.  The retained environments print in this document as Theorem 1 in order of appearance, whereas the article prints the same items with the numbering of its own full document.  The complete native source of the article as distributed in the arXiv e-print for this exact version is bundled unchanged as \texttt{sources/195284/source0.tex}.
\begin{theorem}[$\Delta_0$ bound]\label{thm:Delta0_bound}
For $\pvec,\qvec\in[0,\lam_n]^n$,
\[
\Delta_0\ \le\ \left(2-\frac1n\right)\Delta_1
=\frac{2n-1}{n}\,\Delta_1.
\]
\end{theorem}

\begin{theorem}[$\Delta_2/\Delta_1$ bound]\label{thm:Delta2_bound}
For $\pvec,\qvec\in[0,\lam_n]^n$,
\[
\Delta_2\ \le\ \frac{3(n-1)}{2(2n-1)}\,\Delta_1.
\]
\end{theorem}

\begin{proof}
Fix $a<b$.
Write $\delta_a:=\delta_{\{a\}}$, $\delta_b:=\delta_{\{b\}}$, and $\delta_{ab}:=\delta_{\{a,b\}}$.

We introduce an auxiliary quantity chosen so that a certain linear combination of singleton and doubleton masses isolates $\delta_{ab}$ cleanly in odds coordinates:
\[
S(\yvec;a,b)
:=\beta_n\big(P_{\{a\}}(\yvec)+P_{\{b\}}(\yvec)\big)-2P_{\{a,b\}}(\yvec).
\]
A direct expansion shows
\begin{equation}\label{eq:deltaab_identity}
\frac{\beta_n}{2}(\delta_a+\delta_b)-\delta_{ab}
=\frac12\Big(S(\pvec;a,b)-S(\qvec;a,b)\Big)
\;=:\;\frac12\,\Delta S_{ab}.
\end{equation}
By the triangle inequality,
\begin{equation}\label{eq:deltaab_bound_step1}
\abs{\delta_{ab}}
\le \frac{\beta_n}{2}\big(\abs{\delta_a}+\abs{\delta_b}\big)+\frac12\abs{\Delta S_{ab}}.
\end{equation}
Summing \eqref{eq:deltaab_bound_step1} over all $a<b$ yields
\begin{equation}\label{eq:Delta2_split}
\Delta_2
\le \frac{\beta_n}{2}(n-1)\Delta_1+\sum_{a<b}\frac12\abs{\Delta S_{ab}}.
\end{equation}

It remains to bound $\sum_{a<b}\abs{\Delta S_{ab}}$.
Introduce odds coordinates $o_i(\yvec):=\frac{y_i}{1-y_i}$.
On $[0,\lam_n]$ we have $0\le o_i(\yvec)\le \beta_n$.
Note the factorizations
\[
P_{\{a\}}(\yvec)=P_\emptyset(\yvec)\,o_a(\yvec),\qquad
P_{\{a,b\}}(\yvec)=P_\emptyset(\yvec)\,o_a(\yvec)o_b(\yvec).
\]
Therefore $S(\yvec;a,b)=P_\emptyset(\yvec)\,H(\yvec;a,b)$ where
\[
H(\yvec;a,b):=\beta_n\big(o_a(\yvec)+o_b(\yvec)\big)-2o_a(\yvec)o_b(\yvec).
\]
Write
$\Delta P_\emptyset:=P_\emptyset(\pvec)-P_\emptyset(\qvec)$
and $\Delta o_i:=o_i(\pvec)-o_i(\qvec)$.
Then
\[
\Delta S_{ab}=P_\emptyset(\pvec)\big(H(\pvec;a,b)-H(\qvec;a,b)\big)+H(\qvec;a,b)\,\Delta P_\emptyset.
\]
Using the bilinear identity
\[
H(\pvec;a,b)-H(\qvec;a,b)
=(\beta_n-2o_b(\qvec))\,\Delta o_a+(\beta_n-2o_a(\pvec))\,\Delta o_b,
\]
and the relations
\[
P_\emptyset(\pvec)\,\Delta o_a=\delta_a-o_a(\qvec)\,\Delta P_\emptyset,
\qquad
P_\emptyset(\pvec)\,\Delta o_b=\delta_b-o_b(\qvec)\,\Delta P_\emptyset,
\]
one obtains the identity
\[
\Delta S_{ab}
=(\beta_n-2o_b(\qvec))\delta_a+(\beta_n-2o_a(\pvec))\delta_b+2o_a(\pvec)o_b(\qvec)\,\Delta P_\emptyset.
\]
Taking absolute values and using $o_i(\cdot)\in[0,\beta_n]$ gives
\[
\abs{\Delta S_{ab}}\le \beta_n\abs{\delta_a}+\beta_n\abs{\delta_b}+2\beta_n^2\abs{\Delta P_\emptyset}.
\]
Summing over $a<b$ and using $\abs{\Delta P_\emptyset}=\Delta_0$ yields
\[
\sum_{a<b}\frac12\abs{\Delta S_{ab}}
\le \frac{\beta_n}{2}(n-1)\Delta_1+\beta_n^2\binom{n}{2}\Delta_0.
\]
Apply Theorem~\ref{thm:Delta0_bound} to bound $\Delta_0\le \frac{2n-1}{n}\Delta_1$.
Since $\beta_n=\frac{1}{2n-1}$, we obtain
\[
\beta_n^2\binom{n}{2}\Delta_0
\le \beta_n^2\binom{n}{2}\frac{2n-1}{n}\Delta_1
=\frac{n-1}{2(2n-1)}\Delta_1
=\frac{\beta_n}{2}(n-1)\Delta_1.
\]
Therefore $\sum_{a<b}\frac12\abs{\Delta S_{ab}}\le \beta_n(n-1)\Delta_1$.
Substituting into \eqref{eq:Delta2_split} gives
\[
\Delta_2\le \frac{\beta_n}{2}(n-1)\Delta_1+\beta_n(n-1)\Delta_1
=\frac32\beta_n(n-1)\Delta_1
=\frac{3(n-1)}{2(2n-1)}\Delta_1,
\]
as required.
\end{proof}

\end{document}
